% 第 21 讲 阿廷与布劳尔互反律（上）（源：src/lecture21.txt；已删去抽出时残留的页眉「21. ARTIN AND BRAUER RECIPROCITY, PART I 93/95」）
\section*{第 21 讲\quad 阿廷与布劳尔互反律（上）}

设 $F/\Q$ 是整体域，从而有正合列
\[
1\to F^\times\hookrightarrow A_F^\times\to C_F\to 1,
\]
其中 $C_F:=A_F^\times/F^\times$。上一讲我们几乎证明了：对循环 $n$ 次扩张 $E/F$，有 $\chi(C_E)=n$；剩下的只是证明引理 20.14。

\textbf{引理 21.1.} 设 $L/K$ 是无限域的扩张，$A$ 是 $K$-代数，$M$ 与 $N$ 是 $A$-模，在 $K$ 上有限维，且 $M$ 在 $A$ 上投射。若 $M\otimes_K L\cong N\otimes_K L$ 作为 $A\otimes_K L$-模同构，则 $M\cong N$。

\textbf{证明.} 首先注意有同构
\begin{equation}\tag{21.1}
\operatorname{Hom}_A(M,N)\otimes_K L\xrightarrow{\sim}\operatorname{Hom}_{A\otimes_K L}(M\otimes_K L,\,N\otimes_K L).
\end{equation}
事实上，由于 $M$ 是有限生成投射模，它是某个有限秩自由模的直和项，这就把问题化到 $M=A$ 的情形（因为 Hom 函子与直和交换）。此时 (21.1) 两端都化为 $N\otimes_K L$。（事实上，在更大的一般性下有更基本的恒等式：对任意代数 $A$，只要 $P$ 平坦、$M$ 有限表现，即存在 $n_1,n_2\in\Z$ 使 $A^{n_1}\to A^{n_2}\to M\to 0$ 正合，就有 $\operatorname{Hom}_A(M,N\otimes_A P)=\operatorname{Hom}_A(M,N)\otimes_A P$。）

注意 $M\otimes_K L$ 与 $N\otimes_K L$ 在 $L$ 上有相同的维数，故 $M$ 与 $N$ 在 $K$ 上有相同的维数 $d$。令 $V:=\operatorname{Hom}_A(M,N)$，$W:=\operatorname{Hom}_A(\Lambda^dM,\Lambda^dN)$，其中 $\Lambda^d(-)$ 表示第 $d$ 次外幂。二者都是有限维 $K$-向量空间；特别地，$V$ 是 $d^2$ 维的，$W$ 是 $1$ 维的，即同构于 $K$。$\Lambda^d$ 的函子性给出行列式映射 $\det:V\to W$，它是系数在 $K$ 中的 $d$ 次多项式映射，意思是：取定 $V$ 与 $W$ 的基之后，它由坐标的 $d$ 次多项式给出（因为它是按 $V$ 中 $d\times d$ 矩阵的行列式来计算的）。我们希望证明 $\det$ 非零，因为满足 $\det\varphi\neq0$ 的点 $\varphi\in V$ 恰好就是一个 $A$-模同构 $M\cong N$。由于 $K$ 无限（而 $\det$ 是多项式），只需在将系数扩张到 $L$ 之后检验；由 (21.1)，这给出行列式映射
\[
\operatorname{Hom}_{A\otimes_KL}(M\otimes_KL,\,N\otimes_KL)\to\operatorname{Hom}_{A\otimes_KL}\bigl(\Lambda^d(M\otimes_KL),\,\Lambda^d(N\otimes_KL)\bigr)\cong L.
\]
由于左端含有一个同构 $M\otimes_KL\cong N\otimes_KL$，此映射必非零，证毕。\hfill$\square$

\textbf{证明（引理 20.14）.} 这里的麻烦在于 $\Lambda_1$ 与 $\Lambda_2$ 可能不可公度（即每个都容在另一个的有限指数与倍数之内）。然而我们断言 $\Lambda_1\otimes\Q\cong\Lambda_2\otimes\Q$ 作为 $\Q[G]$-模。事实上，由引理 21.1，只需证明 $\Lambda_1\otimes R\cong\Lambda_2\otimes R$ 作为 $R[G]$-模，而这显然，因为 $\Lambda_i\otimes R\cong V$（$i=1,2$）。于是取 $\Lambda_2$ 在此同构下的像，并与 $R$ 张量以得到 $V$ 中的包含，就存在 $G$-稳定格 $\Lambda_3\subseteq V$，它作为 $\Z[G]$-模同构于 $\Lambda_2$，且与 $\Lambda_1$ 可公度。于是对充分大的 $N$ 有 $N\Lambda_3\subseteq\Lambda_1\subseteq\frac1N\Lambda_3$，故 $\chi(\Lambda_1)=\chi(\Lambda_3)=\chi(\Lambda_2)$，因为所有子商都有限，证毕。\hfill$\square$

现在转向局部 Kronecker--Weber 理论的讨论。

\textbf{定理 21.2.} 对任意素数 $p$，
\[
\Q_p^{\mathrm{ab}}=\Bigl[\bigcup_{n\ge0}\Q_p(\zeta_{p^n})\Bigr]\cdot\Q_p^{\mathrm{unr}},
\]
其中这是与 $\Q_p$ 的未完备化极大非分歧扩张的复合。

\textbf{证明.} 回忆扩张 $\Q_p(\zeta_{p^n})/\Q_p$ 完全分歧，其 Galois 群为 $(\Z/p^n\Z)^\times$。于是 $K\subseteq\Q_p^{\mathrm{ab}}$，其中把右端记为 $K$，且这给出映射
\[
\Z_p^\times\times\widehat\Z=\widehat{\Q}_p^{\times}=\Gal(\Q_p^{\mathrm{ab}}/\Q_p)\longrightarrow\Gal(K/\Q_p)=\Bigl[\varprojlim_n(\Z/p^n\Z)^\times\Bigr]\times\widehat\Z,
\]
其中最左的等号由 lcft 得到，随后是选取 $\Q_p$ 的一个素元；由 Galois 理论该映射满射。由于两端作为抽象群同构，下述引理表明此映射是同构。\hfill$\square$

\textbf{引理 21.3.} 设 $G$ 是投射有限群，使得对一切 $n>0$，$G$ 中指数为 $n$ 的开子群只有有限多个（即 $G$「拓扑有限生成」）。则每个连续同态 $\varphi:G\to G$ 都是同构。

\textbf{证明.} 若 $H\subseteq G$ 是指数至多 $n$ 的子群，则 $\varphi^{-1}(H)\subseteq G$ 也是指数至多 $n$ 的子群。于是有映射
\[
\{H\subseteq G:[G:H]\le n\}\xrightarrow{\ \varphi^{-1}\ }\{H\subseteq G:[G:H]\le n\},
\]
% [原注：原文称此映射「因 $\varphi$ 满射而单射」，但引理假设中 $\varphi$ 只是一般连续同态（如 $\varphi\equiv1$ 则该映射显然不单射）；原文此处有瑕——正确做法是改对 $\varphi(G)$ 施用同一计数论证。译文照译，见报告 §3]
它单射，因为 $\varphi$ 满射。由假设此集合有限，故该映射双射。由于
\[
\operatorname{Im}(\varphi^{-1})=\{H\subseteq G:[G:H]\le n,\ \operatorname{Ker}(\varphi)\subseteq H\},
\]
可知 $\operatorname{Ker}(\varphi)$ 含于 $G$ 的每个有限指数子群，故 $\operatorname{Ker}(\varphi)$ 平凡，$\varphi$ 是同构。\hfill$\square$

现在问：定理 21.2 证明中 $\Z_p^\times\times\widehat\Z$ 的自同构是什么？下述「显式 cft」定理回答此问题，但其证明相当复杂且绝非显然（见 [Dwo58]）。对整体域的类似问题一般未知，只有 $\Q$ 与虚数域的情形除外。

\textbf{定理 21.4（Dwork, Lubin--Tate）.} (1) 元素 $p\in\Q_p^\times$ 在 $\bigcup_n\Q_p(\zeta_{p^n})$ 上平凡作用，而在 $\Q_p^{\mathrm{unr}}$ 上作为 Frobenius 元素作用。(2) 元素 $x\in\Z_p^\times$ 在 $\Q_p^{\mathrm{unr}}$ 上平凡作用，而在 $\bigcup_n\Q_p(\zeta_{p^n})$ 上以 $x^{-1}$ 作用，即 $\theta_p(x)\cdot\zeta_{p^n}=\zeta_{p^n}^{(x^{-1}\bmod p^n)}$，其中
\[
\theta_p:\Q_p^\times\to\Gal\Bigl(\bigcup_n\Q_p(\zeta_{p^n})\cdot\Q_p^{\mathrm{unr}}\Big/\Q_p\Bigr)
\]
是 lcft 给出的同态。

现在我们想介绍两条互反律：阿廷互反律与布劳尔互反律。先从前者开始。设 $E/F$ 是整体域的阿贝尔 $G$-Galois 扩张。回忆我们期待有
\[
F^\times\backslash A_F^\times/N(A_E^\times)=C_F/N(C_E)\xrightarrow{\ \sim\ }G.
\]
我们想构造此映射。

\textbf{断言 21.5.} lcft 给出映射 $\theta:A_F^\times\to G$，满足 $\theta(N(A_E^\times))=1$。

\textbf{证明.} 设 $x\in A_F^\times$。对每个 $v\in M_F$ 有元素 $x_v\in F_v^\times$，lcft 于是给出映射
\[
\theta_v:F_v^\times\to\Gal(E_w/F_v)\subseteq\Gal(E/F)=G
\]
其中 $w\mid v$ 是 $E$ 的一个素位，前者是 $E/F$ 在 $v$ 处的分解群。此嵌入由嵌入 $E\subseteq E_w$ 诱导。回忆当 $E/F$ 阿贝尔时，$\Gal(E_w/F_v)$ 与 $w$ 的选取无关。

我们断言乘积
\[
\theta(x):=\prod_{v\in M_F}\theta_v(x_v)
\]
有意义，即除有限多个 $v$ 外 $\theta_v(x_v)=1$。事实上，对几乎所有的 $v$，$E_w/F_v$ 非分歧且 $x_v\in\mathcal{O}_{F_v}^\times$，从而 $\theta_v(x_v)=1$（因为由 lcft，映射 $\theta_v$ 杀掉 $\mathcal{O}_{F_v}^\times$，并把 $F_v$ 的素元送到 $G$ 的 Frobenius 元素）。由于对一切 $v\in M_F$ 有 $\theta_v(N(E_w^\times))=1$，故 $\theta(N(A_E^\times))=1$，证毕。\hfill$\square$

\textbf{定理 21.6（阿廷互反律）.} 有 $\theta(F^\times)=1$，从而 $\theta$ 给出映射 $C_F/N(C_E)\to G$。

\textbf{例 21.7.} 若 $E/\Q$ 是二次扩张，这归结为二次互反律。事实上，局部阿廷映射恰由希尔伯特符号给出，而我们此前已证明了这一蕴含。

我们将与布劳尔互反律一并证明这一点。
% [原文作 "We will proof this"，疑为 "prove" 之误，见报告 §3]
设 $E/F$ 是整体域的有限 $G$-扩张。有
\[
H^2(G,A_E^\times)=\bigoplus_{v\in M_F}\Br(F_v)=\bigoplus_{v\in M_F\setminus M_F^\infty}\Q/\Z\times\bigoplus_{v\in M_F^\infty}\tfrac12\Z/\Z,
\]
由断言 20.10 以及 $\Br(\R)=\Z/2\Z$。通过
\[
\bigoplus_{v\in M_F/M_F^\infty}\Q/\Z\xrightarrow{\ (x_v)_v\mapsto\sum_vx_v\ }\Q/\Z
\]
定义不变映射 $\iota:H^2(G,A_E^\times)\to\Q/\Z$，即对所有局部因子求和（并忽略一切无限因子）。

\textbf{定理 21.8（布劳尔互反律）.} 复合
\[
\Br(F/E)\to H^2(G,A_E^\times)\xrightarrow{\ \iota\ }\Q/\Z
\]
对一切 $E/\Q$ 为零。
% [原文作 "is zero for all E/Q"；定理设定为 E/F 有限扩张，此处照译，见报告 §3]

注意 $E$ 在这里本质无关紧要；此定理其实关于 $\Br(F)$。第一个映射由对角嵌入 $E\hookrightarrow A_E^\times$ 诱导，而布劳尔群的上同调解释随后表明：$F$ 上每个可除代数除有限多个素位外都是矩阵代数，这本身并不是显然的陈述。

\textbf{例 21.9.} 任意 $a,b\in\Q^\times$ 给出 $\Q$ 上的哈密顿代数 $H_{a,b}$。在有限素位 $p$ 处，不变量为
\[
\begin{cases}
0 & \text{若 } H_{a,b}\otimes\Q_p\cong M_2(\Q_p),\\
1/2 & \text{否则},
\end{cases}
\qquad\Longleftrightarrow\qquad
\begin{cases}
0 & \text{若 } (a,b)_p=1,\\
1/2 & \text{若 } (a,b)_p=-1.
\end{cases}
\]
于是，「所有不变量之和为零」这一断言又是二次互反律。

\textbf{断言 21.10.} 阿廷互反律对 $\Q(\zeta_n)/\Q$ 成立，其中 $\zeta_n$ 是本原 $n$ 次单位根。

\textbf{证明.} 我们借助 Dwork 定理作显式计算。可设 $n=\ell^r$ 是素数幂，因为 $\Q(\zeta_n)$ 是各素数幂因子 $\ell^r\mid n$ 对应的 $\Q(\zeta_{\ell^r})$ 的复合，而 $\Gal(\Q(\zeta_n)/\Q)$ 随后分解为 $\Gal(\Q(\zeta_{\ell^r})/\Q)$ 的积。于是有复合
\[
\Q^\times\hookrightarrow A_\Q^\times\xrightarrow{\ \theta\ }\Gal(\Q(\zeta_{\ell^r})/\Q)=(\Z/\ell^r\Z)^\times,
\]
其中 $\theta=\prod_p\theta_p$ 如前。我们想证明此映射平凡。为此只需证明对一切素数 $p$ 有 $\theta(p)=1$ 且 $\theta(-1)=1$。设 $p\ne\ell$；$p=\ell$ 的情形将在下一讲处理。则
\[
\begin{cases}
\theta_p(p)=p,\\
\theta_\ell(p)=p^{-1},\\
\theta_q(p)=1,\\
\theta_\infty(p)=1.
\end{cases}
\]
事实上，回忆
\[
\theta_p:\Q_p^\times\to\Gal(\Q_p(\zeta_{\ell^r})/\Q_p),\qquad p\mapsto\operatorname{Frob}_p,
\]
且 $\theta_p(\Z_p^\times)=1$。
% [原文作 "θ(Z_p^×) = 1"，疑脱下标 p，见报告 §3]
但 $\operatorname{Frob}_p$ 对应 $p\in(\Z/\ell^r\Z)^\times$。对第二种情形，由 Dwork 定理 $\theta_\ell(p)$ 以 $p^{-1}\in(\Z/\ell^r\Z)^\times$ 作用；对第三种情形，其中 $q\ne p,\ell$，有 $p\in(\Z/q\Z)^\times$ 且扩张 $\Q_q(\zeta_{\ell^r})/\Q_q$ 非分歧。最后，$\theta_\infty$ 对应取符号，因为它是由收缩 $\R^\times$ 的连通分支得到的映射 $\R^\times\to\Gal(\C/\R)$，给出 $\Z/2\Z$ 的、把 $1$ 送到复共轭的映射。\hfill$\square$
